\section{Normalization of data}
	Normalization is the process of scaling numeric values measured in arbitrary units into a standardized relative range, typically between $0$ and $1$ (or $0\%$ and $100\%$).
	
	In optical spectroscopy, X-ray diffraction (XRD), or gamma-ray spectra, raw peak heights depend on lamp power, detector gain, or exposure time. To compare the shapes or relative peak ratios of two different experimental scans, we divide all data values by the maximum peak value:
	\begin{align*}
		Y_\text{normalized}=\dfrac{Y}{Y_\text{max}}
	\end{align*}
	
	When background baseline offset is present, we adjust the scale so the minimum value becomes $0$ and the maximum becomes $1$. This is known as \emph{min-max scaling} as follows:
	\begin{align*}
		Y_\text{normalized}=\dfrac{Y-Y_\text{min}}{Y_\text{max}-Y_\text{min}}
	\end{align*}
	
	\subsection{Normalizing data in Origin}
	
		\subsubsection{First method (in-built method)}
			\begin{enumerate}[I.]
				\item Open the worksheet containing the dataset.
				
				\item Click to highlight the entire column that contains data to be normalized, and do any of the following:
				\begin{enumerate}
					\item From the main menu, go to \verb|Analysis| $>$ \verb|Mathematics| $>$ \verb|Normalize Columns| $>$ \verb|Open Dialog|.
					\item Right-click on the column header and, on the context menu, click on \verb|Normalize|.
				\end{enumerate}
				
				\item Select the appropriate option against \verb|Normalize Methods|.\\ The default is \verb|Normalize to [0, 1]|. Then click \verb|OK|.
				
				\item The normalized data is calculated and placed in a new column. The long name of the column can then be edited as per requirement.
			\end{enumerate}
			
		\subsubsection{Second method (formula method)}
			\begin{enumerate}[I.]
				\item Open the worksheet containing the dataset.
				
				\item Right-click on the empty area to the right of the right-most column, and then click \verb|Add New Column|.
				
				\item A new column with created, say it is named \verb|C|. It can be given an appropriate long name.
				
				\item Click to highlight the new column \verb|C|. Then right-click and select \verb|Set Column Values|.
				
				\item The \verb|Set Values| dialog opens up. Suppose we want to fill up column \verb|C| with normalized values from column \verb|B|, then in the formula box, we type the following formula and click \verb|OK|:
				\begin{verbatim}
		(Col(B) - Min(Col(B))) / (Max(Col(B)) - Min(Col(B)))
				\end{verbatim}
			\end{enumerate}
			
	\subsection*{Exercises}
		\begin{enumerate}
			\item Using \verb|diffraction_scan.csv|, perform normalization on Column B (i.e. \verb|Intensity_counts|) using the Origin's in-built method.
			
			\item Import \verb|muon_events.csv|. Create a new column titled \verb|Normalized_Pulse| and scale \verb|Pulse_Height_mV| using the formula method.
		\end{enumerate}